% =====================================================================
%  1bat-ud5-11.tex  —  SESSIÓ 11: Pèrdua de poder adquisitiu
%  (sense preàmbul: s'inclou des de main.tex)
% =====================================================================
\horatitol{Sessió 11 — Pèrdua de poder adquisitiu}%
{El lloguer puja amb l'IPC però l'ajut es congela: mantenir una xifra «igual» en temps d'inflació és, de fet, perdre poder de compra.}

\subsection{Idea clau}

Recuperem la família de la sessió 10: el lloguer s'ha actualitzat amb l'IPC, però
l'\textbf{ajut social} que rep \textbf{es congela} (no s'actualitza). Avui calculem
què significa això de debò: la \textbf{pèrdua de poder adquisitiu}. És un dels conceptes
més importants i menys visibles de tota la unitat, perquè sovint queda amagat darrere
la frase «l'ajut no ha baixat».

\begin{idea}
\textbf{Poder adquisitiu i inflació}
\begin{itemize}[leftmargin=1.4em, itemsep=2pt]
  \item El \textbf{poder adquisitiu} és el que pots comprar de debò amb uns diners.
  \item Si els preus pugen (IPC positiu) i els teus ingressos es queden \textbf{igual},
        perds poder adquisitiu: amb els mateixos euros compres \emph{menys}.
  \item \textbf{Mantenir una xifra congelada} en temps d'inflació equival, en termes
        reals, a una \textbf{baixada}.
\end{itemize}
\textbf{La frase trampa:} «l'ajut no ha baixat». En euros, no; en poder de compra,
\emph{sí}.
\end{idea}

\subsection{Exemples}

\begin{exemple}[el que li queda a la família]
La família rebia un ajut de $\num{900}\eur$ i pagava $\num{650}\eur$ de lloguer, de
manera que li quedaven $\num{900}-\num{650}=\num{250}\eur$ per a la resta (menjar,
subministraments, transport).

Amb l'IPC, el lloguer puja a $\num{650}\per 1,035=\num{672,75}\eur$, però l'ajut es
manté a $\num{900}\eur$. Ara li queden:
\[
\num{900}-\num{672,75}=\num{227,25}\eur.
\]
Ha perdut $\num{250}-\num{227,25}=\num{22,75}\eur$ \textbf{cada mes} per a tot el que
no és lloguer, és a dir $12\per\num{22,75}=\num{273}\eur$ menys a l'any, \emph{sense
que ningú li hagi «baixat» res}.
\end{exemple}

\begin{exemple}[el valor real d'un ajut congelat]
Mirem-ho des de l'altra banda: quant valen, \emph{en poder de compra}, els
$\num{900}\eur$ de l'ajut un any després, amb una inflació del $3,5\%$? Per saber què
compren \emph{de debò}, els dividim per l'índex de preus:
\[
\frac{\num{900}}{1,035}\approx\num{869,57}\eur.
\]
És a dir, els $\num{900}\eur$ d'ara compren el mateix que $\num{869,57}\eur$ de l'any
passat: l'ajut, tot i ser «el mateix número», ha perdut uns $\num{30,43}\eur$ de poder
de compra.
\end{exemple}

\subsection{Exercicis resolts}

\begin{exresolt}[pèrdua mensual i anual]
Una persona cobra un ajut congelat de $\num{600}\eur$ i paga $\num{400}\eur$ de
lloguer. El lloguer s'actualitza amb un IPC del $4\%$. Quant li quedava abans i quant
li queda ara? Quina és la pèrdua anual?
\solucio
\textbf{Abans:} li quedaven $\num{600}-\num{400}=\num{200}\eur$.

\textbf{Nou lloguer:} $\num{400}\per 1,04=\num{416}\eur$.

\textbf{Ara:} li queden $\num{600}-\num{416}=\num{184}\eur$.

\textbf{Pèrdua:} $\num{200}-\num{184}=\num{16}\eur$ al mes, és a dir
$12\per\num{16}=\num{192}\eur$ a l'any menys per a la resta de despeses.
\resposta{Abans $\num{200}\eur$, ara $\num{184}\eur$; pèrdua de $\num{192}\eur$ l'any.}
\end{exresolt}

\begin{exresolt}[valor real d'una pensió congelada]
Una pensió de $\num{1000}\eur$ es congela un any en què la inflació és del $5\%$. Quant
val, en poder de compra, al cap de l'any? Quant ha perdut?
\solucio
Dividim pel factor d'inflació per saber què compra de debò:
\[
\frac{\num{1000}}{1,05}\approx\num{952,38}\eur.
\]
La pensió «de $\num{1000}\eur$» compra ara el mateix que $\num{952,38}\eur$ de l'any
anterior: ha perdut uns $\num{47,62}\eur$ de poder de compra, tot i no haver baixat ni
un euro en nòmina.
\resposta{Val $\approx\num{952,38}\eur$ en poder de compra; ha perdut uns
$\num{47,62}\eur$.}
\end{exresolt}

\subsection{Exercicis per fer}

\begin{exercicis}
\begin{exlist}
  \item Una família rep un ajut congelat de $\num{800}\eur$ i paga $\num{550}\eur$ de
        lloguer. El lloguer puja amb un IPC del $3\%$. Quant li quedava abans i quant li
        queda ara? Quina és la pèrdua anual?
  \item Una pensió de $\num{1200}\eur$ es congela amb una inflació del $4\%$. Quant val,
        en poder de compra, al cap de l'any? Quant ha perdut?
  \item Un ajut de $\num{500}\eur$ es manté igual durant un any amb un IPC del $6\%$.
        Calcula el seu valor real (dividint per $1,06$) i la pèrdua de poder de compra.
  \item \textbf{(Comparació)} Dos ajuts de $\num{700}\eur$: el primer s'actualitza amb
        l'IPC del $3\%$ i el segon es congela. Quina diferència de poder de compra hi ha
        entre tots dos al cap d'un any?
  \item \textbf{(Interpretació)} Explica, amb les teves paraules, per què «l'ajut no ha
        baixat» pot ser una frase enganyosa en temps d'inflació. Què li ha passat de
        debò a la família?
  \item \textbf{(Raonament — debat)} Com pot una operació aritmètica tan senzilla (un
        percentatge aplicat al lloguer però no a l'ajut) acabar empenyent una família
        cap a l'exclusió social? Per què un tècnic ha d'entendre aquest mecanisme per
        defensar els drets de les persones que acompanya?
\end{exlist}
\end{exercicis}

% ---------------------------------------------------------------------
\begin{idea}
\textbf{Conclusió de l'activitat:} l'aritmètica \textbf{no és neutra}. Un mateix
percentatge aplicat en un lloc i no en un altre descriu —i de vegades determina—
realitats socials. Entendre l'IPC i el poder adquisitiu és una eina per
\textbf{detectar i denunciar} situacions d'injustícia que afecten les famílies
vulnerables.
\end{idea}

% ---------------------------------------------------------------------
\fullsolucions{Sessió 11}
\begin{solucions}
\begin{sollist}
  \item Abans: $\num{800}-\num{550}=\num{250}\eur$. Nou lloguer:
        $\num{550}\per 1,03=\num{566,5}\eur$. Ara: $\num{800}-\num{566,5}=\num{233,5}\eur$.
        Pèrdua: $\num{250}-\num{233,5}=\num{16,5}\eur$ al mes, és a dir
        $12\per\num{16,5}=\num{198}\eur$ l'any.
  \item $\dfrac{\num{1200}}{1,04}\approx\num{1153,85}\eur$ de poder de compra. Ha perdut
        uns $\num{1200}-\num{1153,85}=\num{46,15}\eur$.
  \item Valor real: $\dfrac{\num{500}}{1,06}\approx\num{471,70}\eur$. Pèrdua de poder de
        compra: uns $\num{500}-\num{471,70}=\num{28,30}\eur$.
  \item Ajut actualitzat: $\num{700}\per 1,03=\num{721}\eur$ (manté el poder de compra).
        Ajut congelat: continua a $\num{700}\eur$, que en poder de compra equivalen a
        $\dfrac{\num{700}}{1,03}\approx\num{679,61}\eur$. La diferència de poder de
        compra entre tots dos és d'uns $\num{721}-\num{700}=\num{21}\eur$ nominals al
        mes (o, en termes reals, el congelat ha perdut uns $\num{20,39}\eur$ de poder de
        compra).
  \item Resposta oberta. Idea: en euros l'ajut és el mateix, però com que els preus han
        pujat, amb aquests euros es compra \textbf{menys} que abans; el poder
        adquisitiu ha baixat encara que el número no. De fet, la família té menys
        capacitat real de pagar les seves despeses.
  \item Resposta oberta. Idea: si el lloguer puja amb l'IPC però l'ajut no, cada mes
        queden menys diners per a l'alimentació, els subministraments o el transport;
        acumulat, pot fer que la família no arribi a final de mes i caigui en impagaments
        o privacions. Un tècnic que entén aquest mecanisme pot detectar-ho a temps,
        reclamar l'actualització de l'ajut i defensar els drets de la família.
\end{sollist}
\end{solucions}
